\section{Rung 2: Intervention}

So far, there are several hypergraphs to model a probability distributions by a Bayesian Network.
We now want to understand the directed hyperedges as causal relationships between the variables, beyond encoding of conditional independencies (see \exaref{exa:patientTreatmentOutcome}).
To this end, we model the effects of interventions, which will distinguish between multiple Bayesian Networks modeling the same joint distribution.

\input{examples/patient_treatment_outcome}

Interventions on Bayesian Networks are modeled by the do-calculus \cite{pearl_causality_2009}, which we capture here based on do-variables.

\subsection{Do calculus}

We introduce the do-variables $\dovariableof{\catenumerator}$ taking values in $[\catdimof{\catenumerator}+1]$ for each variable $\catvariableof{\catenumerator}$ of dimension $\catdimof{\catenumerator}$, which model the effects of interventions.
The states are interpreted as
\begin{itemize}
    \item $\doindexof{\catenumerator}\in[\catdimof{\catenumerator}]$: An intervention setting $\catvariableof{\catenumerator}$ to the value $\doindexof{\catenumerator}$
    \item $\doindexof{\catenumerator}=\catdimof{\catenumerator}$: No intervention on $\catvariableof{\catenumerator}$
\end{itemize}

\begin{definition}
    The causally augmented conditional probability core is defined as
    \begin{align*}
        \condprobwrtof{\causalsymbol}{\catvariableof{\atomenumerator}}{\dovariableof{\catenumerator},\catvariableof{\parentsof{\catenumerator}}}
        =& \condprobof{\catvariableof{\catenumerator}}{\catvariableof{\parentsof{\catenumerator}}} \otimes \onehotmapofat{\catdimof{\catenumerator}}{\dovariableof{\catenumerator}} \\
        &+ \sum_{\doindexof{\catenumerator}\in[\catdimof{\catenumerator}]} \onehotmapofat{\doindexof{\catenumerator}}{\catvariableof{\catenumerator}} \otimes \onehotmapofat{\doindexof{\catenumerator}}{\dovariableof{\catenumerator}} \otimes \onesat{\catvariableof{\parentsof{\catenumerator}}}\, .
    \end{align*}
    The causally augmented Bayesian Network is the tensor network
    \begin{align*}
        \contractionof{
            \{\condprobwrtof{\causalsymbol}{\catvariableof{\catenumerator}}{\dovariableof{\catenumerator},\catvariableof{\parentsof{\catenumerator}}} \wcols \catenumeratorin\}
        }{\shortcatvariables,\dovariableof{[\catorder]}} \, .
    \end{align*}
\end{definition}

The Bayesian Networks investigated so far are retrieved by contractions with $\onehotmapofat{\catdimof{\catenumerator}}{\dovariableof{\catenumerator}}$.
This corresponds with the situation, where no interventions are performed on any variable.

\begin{lemma}
    For any Bayesian Network $\probwith$ on a directed acyclic hypergraph $\graph=([\catorder],\edges)$ we have
    \begin{align*}
        \probwith =
        \contractionof{
            \{\condprobwrtof{\causalsymbol}{\catvariableof{\catenumerator}}{\dovariableof{\catenumerator},\catvariableof{\parentsof{\catenumerator}}} \wcols \catenumeratorin\}
            \cup \{\onehotmapofat{\catdimof{\catenumerator}}{\dovariableof{\catenumerator}} \wcols \catenumeratorin \}
        }{\shortcatvariables} \, .
    \end{align*}
\end{lemma}
\begin{proof}
    By construction we have
    \begin{align*}
        \condprobwrtof{\causalsymbol}{\catvariableof{\atomenumerator}}{\dovariableof{\catenumerator}=\catdimof{\catenumerator},\catvariableof{\parentsof{\catenumerator}}}
        =& \condprobof{\catvariableof{\catenumerator}}{\catvariableof{\parentsof{\catenumerator}}}  \, .
    \end{align*}
    Performing this on each causally augmented conditional distribution retrieves thus the original Bayesian Network.
\end{proof}


Adding the do-variables comes with a moderate increase of the bas+ rank.
We provide the following bound.

\begin{lemma}
    For any conditional probability distribution we have
    \begin{align*}
        \bascprankof{\condprobwrtof{\causalsymbol}{\catvariableof{\atomenumerator}}{\dovariableof{\catenumerator},\catvariableof{\parentsof{\catenumerator}}}}
        \leq \condprobat{\catvariableof{\atomenumerator}}{\dovariableof{\catenumerator},\catvariableof{\parentsof{\catenumerator}}}
    \end{align*}
\end{lemma}


\begin{example}[Intervention on a Markov Chain]\label{exa:mcIntervention}

    Let us recall the ambiguity of parameterizing a Markov Chain in \exaref{exa:mcAssociation}, where we observed that both past to future and future to past representations result in Bayesian Networks of similar sparsity.
    Adding do variables to any of these parametrization will result in different interventions (see \figref{fig:mcDoVariableRep}) and thus resolves the ambiguity.

    Intervening on a variable $\catvariableof{\catenumerator}$ and putting it into the state $\catindexof{\catenumerator}$ is done by contracting with
    \begin{align*}
        \onehotmapofat{\catindexof{\catenumerator}}{\dovariableof{\catenumerator}} \, .
    \end{align*}
    We use that since $\catindexof{\catenumerator}\in\catspaceof{\catenumerator}$
    \begin{align*}
        \contractionof{\condprobwrtof{\causalsymbol}{\catvariableof{\catenumerator}}{\catvariableof{\catenumerator-1},\dovariableof{\catenumerator}}, \onehotmapofat{\catindexof{\catenumerator}}{\dovariableof{\catenumerator}} }{
            \catvariableof{\catenumerator},\catvariableof{\catenumerator-1}
        }
        = \onesat{\catvariableof{\catenumerator-1}}
        \otimes  \onehotmapofat{\catindexof{\catenumerator}}{\catvariableof{\catenumerator}}
        %\otimes  \onehotmapofat{\catindexof{\catenumerator}}{\dovariableof{\catenumerator}}
    \end{align*}
    and have that the intervention effectively restarts the Markov Chain at the variable $\catvariableof{\catenumerator}$ (see \figref{fig:mcInterventionExample}a).
    When we choose the causal augmentation of the future-to-past representation instead, the effect would be differently since the restart would be towards the past (see \figref{fig:mcInterventionExample}b).


    \begin{figure}
        \stantikzpic{
            \drawvariablenode{0}{8}{$\dovariableof{\catenumerator-1}$}{DKmin1}
            \drawvariablenode{-3}{5}{$\catvariableof{\catenumerator-1}$}{XKmin1}
            \drawvariablenode{0}{3}{$\dovariableof{\catenumerator}$}{DK}
            \drawvariablenode{-3}{0}{$\catvariableof{\catenumerator}$}{XK}
            \drawvariablenode{0}{-2}{$\dovariableof{\catenumerator+1}$}{DKplus1}
            \drawvariablenode{-3}{-5}{$\catvariableof{\catenumerator+1}$}{XKplus1}

            \draw [->-] (DKmin1)-- (-3,7);
            \draw [->-] (-3,7) -- (XKmin1);
            \draw [->-] (DK)-- (-3,2);
            \draw [->-] (-3,2) -- (XK);
            \draw [->-] (DKplus1) -- (-3,-3);
            \draw [->-] (-3,-3) -- (XKplus1);
            \draw [->-] (XKmin1) -- (-3,2); % (XK);
            \draw [->-] (XK)-- (XKplus1);

            \drawtextnode{-3}{9}{\tiny $\vdots$}
            \draw [->-] (-3,8) -- (-3,7);
            \draw [->-] (XKplus1) -- ++ (0,-3);

            \drawtextnode{-3}{-8.5}{\tiny $\vdots$}

            \begin{scope}[shift={(10,0)}]
                \draw[->-] (-2,10) -- (-2,7) node[midway,right]{\colorlabelsize $\catvariableof{\catenumerator-2}$};
                \draw[->-] (2,8.5) -- (2,7) node[midway,right]{\colorlabelsize $\dovariableof{\catenumerator-1}$};
                \drawtensorblock{0}{6}{4.25}{$\condprobwrtof{\causalsymbol}{\catvariableof{\catenumerator-1}}{\catvariableof{\catenumerator-2},\dovariableof{\catenumerator-1}}$}

                \draw[->-] (-2,5) -- (-2,2) node[midway,right]{\colorlabelsize $\catvariableof{\catenumerator-1}$};
                \draw[->-] (2,3.5) -- (2,2) node[midway,right]{\colorlabelsize $\dovariableof{\catenumerator}$};
                \drawtensorblock{0}{1}{4.25}{$\condprobwrtof{\causalsymbol}{\catvariableof{\catenumerator}}{\catvariableof{\catenumerator-1},\dovariableof{\catenumerator}}$}

                \draw[->-] (-2,0) -- (-2,-3) node[midway,right]{\colorlabelsize $\catvariableof{\catenumerator}$};
                \draw[->-] (2,-1.5) -- (2,-3) node[midway,right]{\colorlabelsize $\dovariableof{\catenumerator+1}$};
                \drawtensorblock{0}{-4}{4.25}{$\condprobwrtof{\causalsymbol}{\catvariableof{\catenumerator+1}}{\catvariableof{\catenumerator},\dovariableof{\catenumerator+1}}$}

                \draw[->-] (-2,-5) -- (-2,-8) node[midway,right]{\colorlabelsize $\catvariableof{\catenumerator+1}$};
            \end{scope}
        }
        \caption{Extension of a Markov Chain in the past-to-future representation by do variables.}\label{fig:mcDoVariableRep}
    \end{figure}


    \begin{figure}[h]
        \stantikzpic{
            \drawtextnode{-4}{9}{$a)$}
            \draw[->-] (-2,10) -- (-2,7) node[midway,right]{\colorlabelsize $\catvariableof{\catenumerator-2}$};
            \draw[->-] (2,8.5) -- (2,7) node[midway,right]{\colorlabelsize $\dovariableof{\catenumerator-1}$};
            \drawtensorblock{0}{6}{4.25}{$\condprobwrtof{\causalsymbol}{\catvariableof{\catenumerator-1}}{\catvariableof{\catenumerator-2},\dovariableof{\catenumerator-1}}$}

            \draw[->-] (-2,5) -- (-2,2) node[midway,right]{\colorlabelsize $\catvariableof{\catenumerator-1}$};
            \draw[->-] (2,3.5) -- (2,2) node[midway,right]{\colorlabelsize $\dovariableof{\catenumerator}$};
            \drawtensorblock{0}{1}{4.25}{$\condprobwrtof{\causalsymbol}{\catvariableof{\catenumerator}}{\catvariableof{\catenumerator-1},\dovariableof{\catenumerator}}$}

            \drawvariabledot{2}{3.5}
            \draw[->-] (5,3.5) -- (2,3.5);
            \drawsquaretensor{6}{3.5}{$\onehotmapof{\catindexof{\catenumerator}}$}

            \draw[->-] (-2,0) -- (-2,-3) node[midway,right]{\colorlabelsize $\catvariableof{\catenumerator}$};

            \draw[->-] (2,-1.5) -- (2,-3) node[midway,right]{\colorlabelsize $\dovariableof{\catenumerator+1}$};
            \drawtensorblock{0}{-4}{4.25}{$\condprobwrtof{\causalsymbol}{\catvariableof{\catenumerator+1}}{\catvariableof{\catenumerator},\dovariableof{\catenumerator+1}}$}

            \draw[->-] (-2,-5) -- (-2,-8) node[midway,right]{\colorlabelsize $\catvariableof{\catenumerator+1}$};

            \drawtextnode{7.5}{1}{${=}$}

            \begin{scope}[shift={(12,0)}]

                \draw[->-] (-2,10) -- (-2,7) node[midway,right]{\colorlabelsize $\catvariableof{\catenumerator-2}$};
                \draw[->-] (2,8.5) -- (2,7) node[midway,right]{\colorlabelsize $\dovariableof{\catenumerator-1}$};


                \drawtensorblock{0}{6}{4.25}{$\condprobwrtof{\causalsymbol}{\catvariableof{\catenumerator-1}}{\catvariableof{\catenumerator-2},\dovariableof{\catenumerator-1}}$}


                \draw[->-] (-2,5) -- (-2,3.5) node[midway,right]{\colorlabelsize $\catvariableof{\catenumerator-1}$};
                \drawsquaretensor{-2}{2.5}{$\ones$}
                \drawsquaretensor{-2}{-0.5}{$\onehotmapof{\catindexof{\catenumerator}}$}
                \draw[->-] (-2,-1.5) -- (-2,-3) node[midway,right]{\colorlabelsize $\catvariableof{\catenumerator}$};

                \draw[->-] (2,-1.5) -- (2,-3) node[midway,right]{\colorlabelsize $\dovariableof{\catenumerator+1}$};
                \drawtensorblock{0}{-4}{4.25}{$\condprobwrtof{\causalsymbol}{\catvariableof{\catenumerator+1}}{\catvariableof{\catenumerator},\dovariableof{\catenumerator+1}}$}

                \draw[->-] (-2,-5) -- (-2,-8) node[midway,right]{\colorlabelsize $\catvariableof{\catenumerator+1}$};

            \end{scope}


            \begin{scope}[shift={(25,0)}]


                \drawtextnode{-4}{9}{$b)$}
                \draw[-<-] (-2,10) -- (-2,7) node[midway,right]{\colorlabelsize $\catvariableof{\catenumerator-2}$};
                \draw[->-] (2,3.5) -- (2,5) node[midway,right]{\colorlabelsize $\dovariableof{\catenumerator-2}$};
                \drawtensorblock{0}{6}{4.25}{$\condprobwrtof{\causalsymbol}{\catvariableof{\catenumerator-2}}{\catvariableof{\catenumerator-1},\dovariableof{\catenumerator-1}}$}

                \draw[-<-] (-2,5) -- (-2,2) node[midway,right]{\colorlabelsize $\catvariableof{\catenumerator-1}$};
                \draw[->-] (2,-1.5) -- (2,0) node[midway,right]{\colorlabelsize $\dovariableof{\catenumerator-1}$};
                \drawtensorblock{0}{1}{4.25}{$\condprobwrtof{\causalsymbol}{\catvariableof{\catenumerator-1}}{\catvariableof{\catenumerator},\dovariableof{\catenumerator}}$}

                \draw[-<-] (-2,0) -- (-2,-3) node[midway,right]{\colorlabelsize $\catvariableof{\catenumerator}$};

                \draw[->-] (2,-6.5) -- (2,-5) node[midway,right]{\colorlabelsize $\dovariableof{\catenumerator}$};
                \drawvariabledot{2}{-6.5}
                \draw[->-] (5,-6.5) -- (2,-6.5);
                \drawsquaretensor{6}{-6.5}{$\onehotmapof{\catindexof{\catenumerator}}$}

                \drawtensorblock{0}{-4}{4.25}{$\condprobwrtof{\causalsymbol}{\catvariableof{\catenumerator}}{\catvariableof{\catenumerator+1},\dovariableof{\catenumerator+1}}$}

                \draw[-<-] (-2,-5) -- (-2,-8) node[midway,right]{\colorlabelsize $\catvariableof{\catenumerator+1}$};

                \drawtextnode{6}{1}{${=}$}

                \begin{scope}[shift={(12,0)}]

                    \draw[-<-] (-2,10) -- (-2,7) node[midway,right]{\colorlabelsize $\catvariableof{\catenumerator-2}$};
                    \draw[->-] (2,3.5) -- (2,5) node[midway,right]{\colorlabelsize $\dovariableof{\catenumerator-2}$};
                    \drawtensorblock{0}{6}{4.25}{$\condprobwrtof{\causalsymbol}{\catvariableof{\catenumerator-2}}{\catvariableof{\catenumerator-1},\dovariableof{\catenumerator-1}}$}

                    \draw[-<-] (-2,5) -- (-2,2) node[midway,right]{\colorlabelsize $\catvariableof{\catenumerator-1}$};
                    \draw[->-] (2,-1.5) -- (2,0) node[midway,right]{\colorlabelsize $\dovariableof{\catenumerator-1}$};
                    \drawtensorblock{0}{1}{4.25}{$\condprobwrtof{\causalsymbol}{\catvariableof{\catenumerator-1}}{\catvariableof{\catenumerator},\dovariableof{\catenumerator}}$}

                    \draw[-<-] (-2,0) -- (-2,-1.5) node[midway,right]{\colorlabelsize $\catvariableof{\catenumerator}$};

                    \drawsquaretensor{-2}{-2.5}{$\onehotmapof{\catindexof{\catenumerator}}$}

                    \drawsquaretensor{-2}{-5.5}{$\ones$}

                    %\drawtensorblock{0}{-4}{4.25}{$\condprobwrtof{\causalsymbol}{\catvariableof{\catenumerator}}{\catvariableof{\catenumerator+1},\dovariableof{\catenumerator+1}}$}

                    \draw[-<-] (-2,-6.5) -- (-2,-8) node[midway,right]{\colorlabelsize $\catvariableof{\catenumerator+1}$};

                \end{scope}

            \end{scope}

        }
        \caption{Effect of intervention on a variable $\dovariableof{\catenumerator}$:
        Restarting the Markov chain at the variable $\catvariableof{\catenumerator}$ in state $\catindexof{\catenumerator}$.
        The effect differs between the past-to-future parameterization a) and the future-to-past parameterization b).
        }
    \end{figure}\label{fig:mcInterventionExample}


\end{example}

\subsection{Intervention Queries}

In a Bayesian Network, when intervening on a variable $\catvariableof{\node}$, the corresponding conditional probability gets trivialized.
The probability tensor is then captured by
\begin{align*}
    \condprobat{\catvariableof{\secnodes}}{\doof{\indexedcatvariableof{\thirdnodes}}}
    %  \condprobat{\catvariableof{\secnodes}}{\doof{X}}
    = \contractionof{
        \{\condprobof{\catvariableof{\node}}{\catvariableof{\parentsof{\node}}} \,:\, \node\notin\thirdnodes \}
        \cup \{\frac{1}{\catdimof{\node}} \onesat{\catvariableof{\node},\catvariableof{\parentsof{\node}}} \,:\, \node\in\thirdnodes\}
        \cup \{\onehotmapofat{\catindexof{\node}}{\catvariableof{\node}} \,:\, \node\in\thirdnodes\}
    }{\catvariableof{\secnodes}} \, .
\end{align*}

Since $\frac{1}{\catdimof{\node}} \onesat{\catvariableof{\node},\catvariableof{\parentsof{\node}}}$ is directed with $\catvariableof{\parentsof{\node}}$ incoming and $\catvariableof{\node}$ outgoing, the partition function of the network stays $1$ and the normalization is captured by the contraction.

%    \subsection{Intervention Variables}
%
%    We modify conditional probability cores in Bayesian Networks to capture interventions, by introducing intervention variable
%
%    \begin{align*}
%        \hypercoreat{\catvariableof{\node},\catvariableof{\parentsof{\node}},\dovariableof{\node}}
%        = \condprobat{\catvariableof{\node}}{\catvariableof{\parentsof{\node}}} \otimes \onehotmapofat{0}{\dovariableof{\node}}
%        + \onesat{\catvariableof{\node},\catvariableof{\parentsof{\node}}} \otimes \tbasisat{\dovariableof{\node}}
%    \end{align*}

\subsection{Simplifications}

In general, when all variables are observable, intervention queries can be expressed by a collection of conditional queries (since any cpd of a Bayesian network is a conditional query, and we contract all except those with interventions on the head variables).
In realistic scenarios, however, not all variables are observable (for example, the smoking gene has been only assumed, but no observations have been made).
We are interested in situations, where intervention queries can be answered by conditional queries of observable variables.
We can proof them in the tensor network formalism based on network separations, which contribution to contractions are scalar multiplications, which therefore are dropped in normalizations.

An example for a simplification is given by the randomized controlled trial.

\begin{example}[Randomized controlled trial]
    The randomized controlled trial is a joint distribution of variables
    \begin{itemize}
        \item Confounding variable $\catvariableof{C}$ summarizing other factors affecting the outcome.
        \item Treatment assignment $\catvariableof{T}$ with values in $\{0,1\}$.
        Here, $1$ indicates assignment to the treatment group and $0$ to the control group.
        The experiment is conducted such that the treatment variable is independent of the confounding variable $\catvariableof{C}$.
        \item Outcome $\catvariableof{O}$ .
    \end{itemize}
    We want to estimate the causal effect of the treatment on the outcome, namely calculate the intervention query
    \begin{align*}
        \condprobat{\catvariableof{O}}{\dovariableof{T}=1} \, .
    \end{align*}
    By design of the randomized control, the confounding variable has no causal effect on $\catvariableof{T}$.
    Thus, we have
    \begin{align*}
        \condprobat{\catvariableof{O}}{\dovariableof{T}=1}
        = \condprobat{\catvariableof{O}}{\catvariableof{T}=1}
    \end{align*}
    that is, the intervention query is equal to the conditional probability (see \figref{fig:RCT}).


    % If confounder influences treatment
    If the treatment variable would depend on the confounder (see \figref{fig:RCTwithConfounder}), the effect of the intervention can not be identified based on conditional probabilities, when not measuring the confounder $\catvariableof{C}$.

    \begin{figure}[hbt!]
        \stantikzpic{

            \drawtextnode{-9}{0}{$a)$}
            \begin{scope}[shift={(-2,-1)}]

                \drawvariablenode{0}{0}{$\catvariableof{O}$}{XO}
                \drawvariablenode{-7}{-5}{$\dovariableof{T}$}{DT}
                \drawvariablenode{-3}{-5}{$\catvariableof{T}$}{XT}
                \drawvariablenode{3}{-5}{$\catvariableof{C}$}{XC}

                \draw[->-] (DT) -- (XT);

                \draw[->-] (XT) -- (0,-2);
                \draw[->-] (0,-2) -- (XO);
                \draw[->-] (XC) -- (0,-2);

            \end{scope}


            \drawtextnode{10}{0}{$b)$}
            \begin{scope}[shift={(17,0)}]
                \drawtdwire{0}{0}{$\catvariableof{0}$}
                \drawtensorblock{0}{-3}{3}{$\condprobat{\catvariableof{O}}{\catvariableof{T},\catvariableof{C}}$}

                \drawtdwire{-2}{-4}{$\catvariableof{T}$}
                \drawtensorblock{-2}{-7}{3}{$\condprobat{\catvariableof{T}}{\dovariableof{T}}$}
                \drawhdwire{-5}{-7}{$\dovariableof{T}$}
                \drawsquaretensor{-8}{-7}{$\onehotmapof{\catindexof{T}}$}

                \drawtdwire{2.5}{-4}{$\catvariableof{C}$}
                \drawtensorblock{3}{-7}{1.5}{$\probat{\catvariableof{C}}$}
%                \drawvariabledot{2}{-6}
            \end{scope}

            \drawtextnode{22}{-5}{${=}$}

            \begin{scope}[shift={(25,0)}]
                \drawtdwire{0}{0}{$\catvariableof{0}$}
                \drawtensorblock{0}{-3}{2.5}{$\condprobat{\catvariableof{O}}{\catvariableof{T}}$}
                \drawtdwire{0}{-4}{$\catvariableof{T}$}
                \drawsquaretensor{0}{-7}{$\onehotmapof{\catindexof{T}}$}
                %\drawtdwire{2}{-4}{$\catvariableof{C}$}
                %\drawvariabledot{2}{-6}
            \end{scope}

        }
        \caption{Model of a Randomized Controlled Trial:
        The confounding variable $\catvariableof{C}$ summarizes all influenced on the outcome $\catvariableof{O}$ other than the treatment variable $\catvariableof{T}$.
        a)
            b) Effect of intervening on the treatement variable $\catindexof{T}\in\{0,1\}$}\label{fig:RCT}
    \end{figure}



    \begin{figure}[hbt!]
        \stantikzpic{

            \drawtextnode{-9}{0}{$a)$}
            \begin{scope}[shift={(-2,-1)}]

                \drawvariablenode{0}{0}{$\catvariableof{O}$}{XO}
                \drawvariablenode{-7}{-6}{$\dovariableof{T}$}{DT}
                \drawvariablenode{-3}{-3}{$\catvariableof{T}$}{XT}
                \drawvariablenode{1}{-6}{$\catvariableof{C}$}{XC}

                \draw[->-] (DT) -- (-3,-5);
                \draw[->-] (-3,-5) -- (XT);
                \draw[->-] (XC) -- (-3,-5);

                \draw[->-] (XT) -- (0,-2);
                \draw[->-] (0,-2) -- (XO);
                \draw[->-] (XC) -- (0,-2);

            \end{scope}


            \drawtextnode{10}{0}{$b)$}
            \begin{scope}[shift={(17,0)}]
                \drawtdwire{0}{0}{$\catvariableof{0}$}
                \drawtensorblock{0}{-3}{3}{$\condprobat{\catvariableof{O}}{\catvariableof{T},\catvariableof{C}}$}

                \drawtdwire{-2}{-4}{$\catvariableof{T}$}
                \drawtensorblock{-2}{-7}{3}{$\condprobat{\catvariableof{T}}{\catvariableof{C},\dovariableof{T}}$}
                \drawhdwire{-5}{-7}{$\dovariableof{T}$}
                \drawsquaretensor{-8}{-7}{$\onehotmapof{\catindexof{T}}$}

                \drawtdwire{2}{-4}{$\catvariableof{C}$}
                \draw  (2,-10) -- (2,-6);
                \drawvariabledot{2}{-10}
                \draw [->-] (2,-11) -- (2,-10);
                \draw [->-] (2,-10) to [bend right=-30] (-2,-8);
                \drawtensorblock{2}{-12}{1.5}{$\probat{\catvariableof{C}}$}

            \end{scope}

            \drawtextnode{22}{-5}{${=}$}

            \begin{scope}[shift={(27,0)}]
                \drawtdwire{0}{0}{$\catvariableof{0}$}
                \drawtensorblock{0}{-3}{3}{$\condprobat{\catvariableof{O}}{\catvariableof{T},\catvariableof{C}}$}
                \drawtdwire{-2}{-4}{$\catvariableof{T}$}
                \drawsquaretensor{-2}{-7}{$\onehotmapof{\catindexof{T}}$}
                \drawtdwire{2}{-4}{$\catvariableof{C}$}
                %\drawvariabledot{2}{-6}
                %\drawtdwire{2.5}{-4}{$\catvariableof{C}$}
                \drawtensorblock{2}{-7}{1.5}{$\probat{\catvariableof{C}}$}
            \end{scope}

        }
        \caption{Effect when the confounder has a causal influence on the treatment:
        The causal effect is not identifiable when not measuring the confounding variable.}\label{fig:RCTwithConfounder}
    \end{figure}

\end{example}

\subsubsection{Backdoor Criterion}

More general than in the randomized controlled trial we have to deal with confounding variables.
The Backdoor Criterion and the Frontdoor Criterion exploit mediator variables $\catvariableof{M}$, which are also observable and block certain causal paths to unobserved confounding variables $\catvariableof{U}$.

\begin{definition}\label{def:backdoorCriterion}
    Let there be a directed acyclic hypergraph, such that $\catvariableof{0}$ is topologically ordered before $\catvariableof{1}$.
    A backdoor path is a path through the hypergraph, such that the first and the second node in the path are in the outgoing and the incoming nodes of a hyperedge.
    If a variable from a set of variables $\catvariableof{M}$ is contained in every backdoor path from $\catvariableof{0}$ to $\catvariableof{1}$, then $\catvariableof{M}$ satisfies the backdoor criterion with respect to $\catvariableof{0}$ and $\catvariableof{1}$.
\end{definition}

See \figref{fig:backdoorCriterion} for a typical hypergraph where $\catvariableof{M}$ satisfies the backdoor criterion.

\begin{lemma}\label{lem:bdCriterionMediatorCond}
    When $\catvariableof{M}$ satisfies the backdoor criterion for $\catvariableof{0}$ and $\catvariableof{1}$, and for some $\catindexofin{M},\catindexofin{0}$ we have $\probat{\indexedcatvariableof{0},\indexedcatvariableof{M}\neq 0}$ then
    \begin{align*}
        \condprobat{\catvariableof{1}}{\indexedcatvariableof{0},\indexedcatvariableof{M}}
        = \condprobwrtof{\causalsymbol}{\catvariableof{1}}{\dovariableof{0}=\catindexof{0},\indexedcatvariableof{M}} \, .
    \end{align*}
%    and
%    \begin{align*}
%        \condprobat{\catvariableof{M}}{\dovariableof{0}=\catindexof{0}}
%        = \probat{\catvariableof{M}} \, .
%    \end{align*}
\end{lemma}
\begin{proof}
        %To show the first assertions,
    We compare the contractions $\probat{\catvariableof{1},\indexedcatvariableof{0},\indexedcatvariableof{M}}$ and  $\probofat{\causalsymbol}{\catvariableof{1},\dovariableof{0}=\catindexof{0},\indexedcatvariableof{M}}$, and notice that both differ only in a factor.
    Whenever $\probat{\indexedcatvariableof{0},\indexedcatvariableof{M}}\neq 0$ the factor does not vanish.
    After normalization with $\catvariableof{0}$ and $\catvariableof{M}$ incoming, the tensors are therefore equivalent.
        %We notice that adding $\onehotmapofat{\catindexof{M}}{\catvariableof{M}}$ decomposes the contraction into a contraction staying the same fo


\end{proof}


\begin{theorem}[Backdoor Criterion]
    Let $\catvariableof{M}$ satisfy the backdoor criterion with respect to $\catvariableof{0}$ and $\catvariableof{1}$, no node in $\catvariableof{M}$ be a descendant of $\catvariableof{0}$, and let $\catindexofin{0}$ be such that for any $\catindexofin{M}$ we have $\probat{\indexedcatvariableof{0},\indexedcatvariableof{M}}\neq0$.
    Then we have
    %blocks any backdoor path from $\catvariableof{0}$ to $\catvariableof{1}$ and no node in $\catvariableof{M}$ is a descendant of $\catvariableof{0}$, then
    \begin{align*}
        \condprobwrtof{\causalsymbol}{\catvariableof{1}}{\dovariableof{0}=\catindexof{0}}
        = \contractionof{
            \condprobof{\catvariableof{1}}{\indexedcatvariableof{0},\catvariableof{M}},
            \probat{\catvariableof{M}}
        }{\catvariableof{1}} \, .
    \end{align*}
\end{theorem}
\begin{proof}
    We have by the Bayes rule
    \begin{align*}
        \condprobwrtof{\causalsymbol}{\catvariableof{1}}{\dovariableof{0}=\catindexof{0}}
        = \contractionof{
            \condprobwrtof{\causalsymbol}{\catvariableof{1}}{\catvariableof{M},\dovariableof{0}=\catindexof{0}},
            \condprobwrtof{\causalsymbol}{\catvariableof{M}}{\dovariableof{0}=\catindexof{0}}
        }{\catvariableof{1}} \, .
    \end{align*}
    We apply \lemref{lem:bdCriterionMediatorCond} for arbitrary $\catindexofin{M}$ and get
    \begin{align*}
        \condprobwrtof{\causalsymbol}{\catvariableof{1}}{\indexedcatvariableof{M},\dovariableof{0}=\catindexof{0}}
        = \condprobat{\catvariableof{1}}{\indexedcatvariableof{M},\indexeddovariableof{0}} \, .
    \end{align*}
    Further, we notice that the contraction $\condprobwrtof{\causalsymbol}{\catvariableof{M}}{\dovariableof{0}=\catindexof{0}}$ simplifies to $\probat{\catvariableof{M}}$ given the assumptions (by exploiting the directed notation to iteratively trivialize the conditional probability tensors).


%    By assumption, we can contract all variables beside $\catvariableof{[2]},\catvariableof{M}$ into cpds, and have a simplified causal model (see \figref{fig:backdoorCriterion}a).
%    We have
%    \begin{align*}
%        \condprobwrtof{\causalsymbol}{\catvariableof{0}}{\dovariableof{0}=\catindexof{0},\catvariableof{M}}
%        = \onehotmapofat{\catindexof{0}}{\catvariableof{0}} \otimes \onesat{\catvariableof{M}}
%    \end{align*}
%    and thus
%    \begin{align*}
%        \condprobwrtof{\causalsymbol}{\catvariableof{1}}{\dovariableof{0}=\catindexof{0}}
%        = \contractionof{
%            \condprobof{\catvariableof{1}}{\indexedcatvariableof{0},\catvariableof{M}},
%            \probat{\catvariableof{M}}
%        }{\catvariableof{1}} \, .
%    \end{align*}
\end{proof}

\begin{figure}[t]
    \stantikzpic{
    \node[anchor=center] at (-7,3) {$a)$};

    \drawvariablenode{-3}{-3}{$\dovariableof{\intervenedset}$}{D0}

    \drawvariablenode{0}{0}{$\intervenedvar$}{X0}
    \coordinate (E0) at ($(X0) - (0,2)$);

    \drawvariablenode{4}{3}{$\observedvar$}{X1}
    \coordinate (E1) at ($(X1) - (0,2)$);

    \drawvariablenode{3}{-3}{$\mediatorvar$}{M}
    \drawvariablenode{8}{0}{$\confoundervar$}{U}

    \draw[->-] (D0) -- (E0);
    \draw[->-] (M) -- (E0);
    \draw[->-] (E0) -- (X0);

    \draw[->-] (X0) -- (E1);
    \draw[->-] (U) -- (E1);
    \draw[->-] (E1) -- (X1);

    \draw[->-] (M) -- (U);

    \begin{scope}[shift={(20,0)}]

        \node[anchor=center] at (-2,3) {$b)$};

        \drawvariablenode{0}{0}{$\intervenedvar$}{X0}
        \coordinate (E0) at ($(X0) - (0,2)$);

        \drawvariablenode{4}{3}{$\observedvar$}{X1}
        \coordinate (E1) at ($(X1) - (0,2)$);

        \drawvariablenode{3}{-3}{$\mediatorvar$}{M}
        \drawvariablenode{8}{0}{$\confoundervar$}{U}

        \draw[->-] (X0) -- (E1);
        \draw[->-] (U) -- (E1);
        \draw[->-] (E1) -- (X1);

        \draw[->-] (M) -- (U);
    \end{scope}
}


    \caption{Causal graph for the Backdoor Criterion (see \defref{def:backdoorCriterion}).
    a) The mediator variable $\catvariableof{M}$ blocks any backdoor path from $\catvariableof{0}$ to $\catvariableof{1}$, potentially through a confounding variable $\catvariableof{U}$.
    b) The effect of the intervention on $\catvariableof{0}$ reduces the Bayesian network effectively to the shown subgraph.}
    \label{fig:backdoorCriterion}
\end{figure}

\subsubsection{Frontdoor Criterion}

\begin{definition}
    A frontdoor path from $\catvariableof{0}$ to $\catvariableof{1}$ is any directed path from $\catvariableof{0}$ to $\catvariableof{1}$.
    A set $\catvariableof{M}$ of variables satisfies the frontdoor criterion, if any frontdoor path from $\catvariableof{0}$ to $\catvariableof{1}$ contains a node in $\catvariableof{M}$.
\end{definition}


\begin{lemma}[Frontdoor Criterion]
    When $\catvariableof{M}$ satisfies the frontdoor criterion with respect to $\catvariableof{0}$ and $\catvariableof{1}$, and there is no confounder between $\catvariableof{0}$ and $\catvariableof{M}$, as well as between $\catvariableof{M}$ and $\catvariableof{1}$, then
    \begin{align*}
        \condprobwrtof{\causalsymbol}{\catvariableof{1}}{\dovariableof{0}=\catindexof{0}}
        = \contractionof{
            \condprobof{\catvariableof{M}}{\indexedcatvariableof{0}},
            \contractionof{\condprobof{\catvariableof{1}}{\catvariableof{M},\catvariableof{0}},\probat{\catvariableof{0}}}{\catvariableof{1},\catvariableof{M}}
%            \condprobof{\catvariableof{1}}{\catvariableof{M},\tildecatvariableof{0}},
%            \probat{\tildecatvariableof{0}}
        }{\catvariableof{1}} \, .
    \end{align*}
    Note that the inner contraction has $\catvariableof{0}$ as a closed variable.
\end{lemma}
\begin{proof}
    We show the claim for the situation in \figref{fig:frontdoorCriterion}
    \begin{align*}
        \condprobwrtof{\causalsymbol}{\catvariableof{1}}{\dovariableof{0}=\catindexof{0}}
        &= \contractionof{
            \condprobof{\catvariableof{M}}{\indexedcatvariableof{0}},
            \condprobof{\catvariableof{1}}{\catvariableof{M},\catvariableof{U}},
            \probat{\catvariableof{U}}
        }{\catvariableof{1}} \\
        &= \contractionof{
            \condprobof{\catvariableof{M}}{\indexedcatvariableof{0}},
            \contractionof{\condprobof{\catvariableof{1}}{\catvariableof{M},\catvariableof{U}},\probat{\catvariableof{U}}}{\catvariableof{1},\catvariableof{M}}
        }{\catvariableof{1}} \, .
    \end{align*}
    We note that
    \begin{align*}
        \contractionof{\condprobof{\catvariableof{1}}{\catvariableof{M},\catvariableof{U}},\probat{\catvariableof{U}}}{\catvariableof{1},\indexedcatvariableof{M}}
        &= \condprobwrtof{\causalsymbol}{\catvariableof{1}}{\dovariableof{M}=\catindexof{M}} \\
        &= \contractionof{\condprobof{\catvariableof{1}}{\catvariableof{M},\catvariableof{0}},\probat{\catvariableof{0}}}{\catvariableof{1},\indexedcatvariableof{M}}
    \end{align*}
    Here we used in the second equation that the variable $\catvariableof{0}$ blocks all backdoor paths from $\catvariableof{M}$ to $\catvariableof{1}$ (i.e. the backdoor criterion).
    Combining both equations proves the claim.
\end{proof}


\begin{figure}[t]
    \stantikzpic{
        \drawvariablenode{-3}{-3}{$\dovariableof{0}$}{D0}

        \drawvariablenode{0}{0}{$\catvariableof{0}$}{X0}
        \coordinate (E0) at ($(X0) - (0,2)$);

        \drawvariablenode{10}{4}{$\catvariableof{1}$}{X1}
        \coordinate (E1) at ($(X1) - (1,2)$);

        \drawvariablenode{5}{2}{$\catvariableof{M}$}{M}
        \drawvariablenode{5}{-3}{$\catvariableof{U}$}{U}

        \draw[->-] (D0) -- (E0);
        \draw[->-] (U) -- (E0);
        \draw[->-] (E0) -- (X0);

        \draw[->-] (X0) -- (M);

        \draw[->-] (M) -- (E1);
        \draw[->-] (U) -- (E1);
        \draw[->-] (E1) -- (X1);
}
    \caption{Typical hypergraph, where $\catvariableof{M}$ satisfies the frontdoor criterion with respect to $\catvariableof{0}$ and $\catvariableof{1}$.}
    \label{fig:frontdoorCriterion}
\end{figure}